% Experiment 1: UNIX Permission and umask Calculator

\section{Experiment 1: UNIX Permission and umask Calculator}

\subsection{Problem Statement}
The objective is to implement a CLI tool named \texttt{permcalc} that simulates how UNIX permissions are calculated when applying a \texttt{umask}.

The tool must:
\begin{itemize}
    \item Accept inputs: \verb|--mode <octal>| (required) and optional \verb|--umask <octal>|.
    \item Validate that \verb|<octal>| is exactly 4 digits from 0000 to 0777.
    \item Compute the effective permission: \verb|effective_mode = mode & (~umask)|.
    \item Output two lines on success: \verb|OK: EFFECTIVE <octal>| and \verb|OK: SYMBOLIC <rwxrwxrwx>|.
    \item Handle errors using specific codes: \verb|E_USAGE|, \verb|E_OCTAL|, \verb|E_RANGE|.
\end{itemize}

\subsection{Source Code}
\begin{lstlisting}[language=C, caption={permcalc.c}]
// #####  Permission Calculation  ######

#include <stdio.h>
#include <stdlib.h>
#include <string.h>
#include <ctype.h>

// Error codes 
#define E_USAGE "E_USAGE"
#define E_OCTAL "E_OCTAL"
#define E_RANGE "E_RANGE"

// Error Printing
void print_error(const char *code, const char *msg) {
    fprintf(stderr, "ERROR: %s: %s\n", code, msg);
    exit(1);
}

// Validates 4 digit octal
void validate_octal(const char *str) {
    if (strlen(str) != 4) {
        print_error(E_OCTAL, "mode must be 4-digit octal (0000-0777)");
    }
    for (int i = 0; i < 4; i++) {
        if (str[i] < '0' || str[i] > '7') {
            print_error(E_OCTAL, "mode must be 4-digit octal (0000-0777)");
        }
    }
}

// Converting Octal String to Int
int parse_octal(const char *str) {
    validate_octal(str);
    int val = strtol(str, NULL, 8);
    if (val < 0 || val > 0777) {
        print_error(E_RANGE, "value out of range 0000-0777");
    }
    return val;
}

// Generating rwxrwxrwx string
void get_symbolic(int mode, char *buffer) {
    const char *perms = "rwx";
    for (int i = 0; i < 9; i++) {
        if (mode & (1 << (8 - i))) {
            buffer[i] = perms[i % 3];
        } else {
            buffer[i] = '-';
        }
    }
    buffer[9] = '\0';
}

int main(int argc, char *argv[]) {
    char *mode_str = NULL;
    char *umask_str = NULL;

    // Parsing CLI arguments
    for (int i = 1; i < argc; i++) {
        if (strcmp(argv[i], "--mode") == 0) {
            if (i + 1 < argc) mode_str = argv[++i];
            else print_error(E_USAGE, "missing value for --mode");
        } else if (strcmp(argv[i], "--umask") == 0) {
            if (i + 1 < argc) umask_str = argv[++i];
            else print_error(E_USAGE, "missing value for --umask");
        } else {
            // Strict usage checking
             print_error(E_USAGE, "unknown argument");
        }
    }

    if (!mode_str) {
        print_error(E_USAGE, "missing required --mode");
    }

    int mode = parse_octal(mode_str);
    int umask = 0000; // Taking this value as default
    if (umask_str) {
        umask = parse_octal(umask_str);
    }
    // Masking with 0777 to avoid integer sign extension issues
    int effective = (mode & (~umask)) & 0777;

    char symbolic[10];
    get_symbolic(effective, symbolic);

    // Output success message
    printf("OK: EFFECTIVE %04o\n", effective);
    printf("OK: SYMBOLIC %s\n", symbolic);

    return 0;
}
\end{lstlisting}

\subsection{Sample Input and Output}

\textbf{Case 1}
\begin{verbatim}
$ ./exp_01 --mode 0644 --umask 002 
OK: EFFECTIVE 0644
OK: SYMBOLIC rw-r--r--
\end{verbatim}

\textbf{Case 2}
\begin{verbatim}
$ ./exp_01 --mode 644 --umask 002 
ERROR: E_OCTAL: mode must be 4-digit octal (0000-0777)
\end{verbatim}